\documentclass[letterpaper,11pt]{article}

\usepackage{latexsym}
\usepackage[empty]{fullpage}
\usepackage{titlesec}
\usepackage{marvosym}
\usepackage[usenames,dvipsnames]{color}
\usepackage{verbatim}
\usepackage{enumitem}
\usepackage[hidelinks]{hyperref}
\usepackage{fancyhdr}
\usepackage[english]{babel}
\usepackage{tabularx}
\usepackage{fontawesome5}
\usepackage{multicol}
\setlength{\multicolsep}{-3.0pt}
\setlength{\columnsep}{-1pt}
\input{glyphtounicode}


%----------FONT OPTIONS----------
% sans-serif
% \usepackage[sfdefault]{FiraSans}
% \usepackage[sfdefault]{roboto}
% \usepackage[sfdefault]{noto-sans}
% \usepackage[default]{sourcesanspro}

% serif
% \usepackage{CormorantGaramond}
% \usepackage{charter}


\pagestyle{fancy}
\fancyhf{} % clear all header and footer fields
\fancyfoot{}
\renewcommand{\headrulewidth}{0pt}
\renewcommand{\footrulewidth}{0pt}

% Adjust margins
\addtolength{\oddsidemargin}{-0.6in}
\addtolength{\evensidemargin}{-0.5in}
\addtolength{\textwidth}{1.19in}
\addtolength{\topmargin}{-.7in}
\addtolength{\textheight}{1.4in}

\urlstyle{same}

\raggedbottom
\raggedright
\setlength{\tabcolsep}{0in}

% Sections formatting
\titleformat{\section}{
  \vspace{-4pt}\scshape\raggedright\large\bfseries
}{}{0em}{}[\color{black}\titlerule \vspace{-5pt}]

% Ensure that generate pdf is machine readable/ATS parsable
\pdfgentounicode=1

%-------------------------
% Custom commands
\newcommand{\resumeItem}[1]{
  \item\small{
    {#1 \vspace{-2pt}}
  }
}

\newcommand{\classesList}[4]{
    \item\small{
        {#1 #2 #3 #4 \vspace{-2pt}}
  }
}

\newcommand{\resumeSubheading}[4]{
  \vspace{-2pt}\item
    \begin{tabular*}{1.0\textwidth}[t]{l@{\extracolsep{\fill}}r}
      \textbf{#1} & \textbf{\small #2} \\
      \textit{\small#3} & \textit{\small #4} \\
    \end{tabular*}\vspace{-7pt}
}

\newcommand{\resumeSubSubheading}[2]{
    \item
    \begin{tabular*}{0.97\textwidth}{l@{\extracolsep{\fill}}r}
      \textit{\small#1} & \textit{\small #2} \\
    \end{tabular*}\vspace{-7pt}
}

\newcommand{\resumeProjectHeading}[2]{
    \item
    \begin{tabular*}{1.001\textwidth}{l@{\extracolsep{\fill}}r}
      \small#1 & \textbf{\small #2}\\
    \end{tabular*}\vspace{-7pt}
}

\newcommand{\resumeSubItem}[1]{\resumeItem{#1}\vspace{-4pt}}

\renewcommand\labelitemi{$\vcenter{\hbox{\tiny$\bullet$}}$}
\renewcommand\labelitemii{$\vcenter{\hbox{\tiny$\bullet$}}$}

\newcommand{\resumeSubHeadingListStart}{\begin{itemize}[leftmargin=0.0in, label={}]}
\newcommand{\resumeSubHeadingListEnd}{\end{itemize}}
\newcommand{\resumeItemListStart}{\begin{itemize}}
\newcommand{\resumeItemListEnd}{\end{itemize}\vspace{-5pt}}

%-------------------------------------------
%%%%%%  RESUME STARTS HERE  %%%%%%%%%%%%%%%%%%%%%%%%%%%%


\begin{document}

%----------HEADING----------
% \begin{tabular*}{\textwidth}{l@{\extracolsep{\fill}}r}
%   \textbf{\href{http://sourabhbajaj.com/}{\Large Sourabh Bajaj}} & Email : \href{mailto:sourabh@sourabhbajaj.com}{sourabh@sourabhbajaj.com}\\
%   \href{http://sourabhbajaj.com/}{http://www.sourabhbajaj.com} & Mobile : +1-123-456-7890 \\
% \end{tabular*}

\begin{center}
    {\Huge \scshape Daniel Duclos-Cavalcanti} \\ \vspace{1pt}
    \textbf{\large{\scshape Systems Software Engineer}} \\ \vspace{1pt}
    New York, New York \\ \vspace{1pt}
    \small 
    % XXX-XXX-XXXX 
     516-912-7975 
    $|$ U.S. Citizen $|$
    \href{mailto:daniel@duclos.dev}{\underline{daniel@duclos.dev}} $|$ 
    \href{https://www.duclos.dev}{\underline{www.duclos.dev}} $|$
    \href{https://www.linkedin.com/in/duclos-cavalcanti/}{\underline{linkedin/duclos-cavalcanti}} $|$
    \href{https://github.com/duclos-cavalcanti}{\underline{github/duclos-cavalcanti}} 
\end{center}

% \section{Summary}
% \small{
% Creative thinker and problem-solver with a masters and bachelors in 
% computer engineering from Germany. Today, I am in New York, 
% collaborating on research with Dr.Sivaraman (NYU) on distributed 
% low-latency networking on the cloud.
% }

\vspace{-16pt}
%-----------EXPERIENCE-----------
\vspace{-4pt}
\section{Experience}
    \resumeSubHeadingListStart
        \resumeSubheading
            {Software Engineer II}{Mar 2025 -- Present}
            {Ripple}{New York, USA}
                \resumeSubSubheading
                    {Managed Custody}{Mar 2026 -- Present}
                \resumeItemListStart
                \resumeItem{One of \textbf{3} engineers building a complete \textbf{greenfield} managed-custody product: a \textbf{Rust} orchestration layer over Ripple Custody on \textbf{AWS Lambda}, \textbf{Step Functions}, \textbf{EventBridge}, \textbf{DynamoDB}, and \textbf{KMS} for tier-1/2 bank clients.}
                \resumeItem{Owned the \textbf{actors}, \textbf{approvals}, \textbf{events}, and \textbf{domains} APIs: \textbf{casbin}-gated maker/checker \textbf{quorums}, per-actor \textbf{KMS} signing keys bound by \textbf{proof-of-possession}, and an \textbf{event-sourced}, cursor-paginated stream.}
                \resumeItem{Built and managed scalable infrastructure using \textbf{AWS CDK}: \textbf{35 stacks}, $\sim$1,100 resources per environment, branch-scoped naming allowing highly concurrent \textbf{ephemeral environment} deployments with auto-tear-down.}
                \resumeItem{Architected an \href{https://github.com/duclos-cavalcanti/agent-factory}{\underline{AI-native software factory}}: a crew of role-specialized agents (\textbf{plan $\to$ implement $\to$ review $\to$ merge}) on a cron-driven state machine semi-autonomously shipping \textbf{Jira} tickets to production, lifting delivery \textbf{$\sim$10$\times$} (\textbf{15}$\to$\textbf{155 MRs/mo}, peak \textbf{205}) across the product's \textbf{$\sim$6-month} build.}
                \resumeItemListEnd

                \resumeSubSubheading
                    {Ripple Custody (Core Platform)}{Mar 2025 -- Mar 2026}
                \resumeItemListStart
                \resumeItem{Migrated a \textbf{HSM}-backed key \textbf{cold-storage} custody service (vault) from legacy \textbf{C\#/.NET} to a \textbf{Rust} (\texttt{axum}) drop-in replacement, serving its batch file-transfer execution API over \textbf{gRPC} to two live customers.}
                \resumeItem{Developed a benchmarking harness in \textbf{Rust} (\texttt{criterion}) to profile cryptographic operations across the custody's \textbf{KMS} backends (\textbf{IBM} and \textbf{Luna} \textbf{HSMs}, \textbf{MPC}), tracking performance regressions in \textbf{CI}.}
                \resumeItemListEnd

        \resumeSubheading
            {Research Assistant}{Jul 2022 -- Oct 2022}
            {TU Munich}{Munich, Germany}
            \resumeItemListStart
                \resumeItem{Profiled ML inference across CPUs, GPUs, and Coral Edge TPUs (TensorFlow Lite) via USB traffic analysis (PyShark).}
                \resumeItem{Established up to 57\% of inference time is data transmission to the TPU, informing TensorDSE's hardware mapping.}
                % \resumeItem{Collaborated on \href{https://github.com/alxhoff/TensorDSE}{\underline{TensorDSE}}, a Design-Space Exploration framework to accelerate machine learning model deployments.}
                % \resumeItem{Assessed the performance metrics of multiple ML models across GPUs, CPUs and TPUs with TensorFlow Lite.}
                % \resumeItem{Generated inference cost analysis reports on Google's Coral Edge TPU via USB traffic analysis (PyShark).}
                % \resumeItem{Established that on average up to 57\% of total inference time consists of data transmission with the external TPU.}
                % \resumeItem{Generated cost analysis reports on Google's Coral Edge TPU via USB traffic analysis (PyShark) during transmission.}
                % \resumeItem{Discovered that over 60\% of total inference time was attributed to data transmission with the external TPU.}
                % \resumeItem{TensorDSE consumed reports to map a model's deployment optimally onto an available set of hardware devices.}
            \resumeItemListEnd

        \resumeSubheading
            {Embedded Software Engineer -- Part Time}{Aug 2021 -- Jan 2022}
            {Molabo GmbH}{Ottobrunn, Germany}
            \resumeItemListStart
                \resumeItem{Built simulation tooling via Linux virtual CAN, raising GTest coverage 25\% on safety-critical motor-controller firmware.}
                \resumeItem{Owned the team's Jenkins/Docker/CMake \textbf{CI} pipeline, supporting 10+ engineers.}
                % \resumeItem{Extended the \textbf{CAN-bus} firmware update system for 18+ clients, hardening partial-update reliability.}
                % \resumeItem{Extended the firmware update system for over 18 clients, enhancing reliability of partial updates via CAN bus.}
                % \resumeItem{Developed robust firmware update systems, ensuring reliability for partial updates (CAN bus) across 18+ clients.}
                % \resumeItem{Actively developed the team's CI/CD build pipeline via Jenkins, Docker, and CMake, supporting over 10+ engineers.}
            \resumeItemListEnd

        \resumeSubheading
            {Tutor (Embedded Systems Programming Lab)}{Apr 2021 -- Aug 2021}
            {TU Munich}{Munich, Germany}
            \resumeItemListStart
                % \resumeItem{Mentored over 12 students on designing and developing low-level embedded FreeRTOS applications in C.}
                % \resumeItem{Conducted 30+ sessions on best practices in data structures, concurrency, performance, and real-time scheduling.}
                \resumeItem{Taught 30+ lab sessions to 12+ students on real-time, multi-threaded FreeRTOS development in C.}
            \resumeItemListEnd
    \resumeSubHeadingListEnd 

\vspace{-14pt}
\section{Publications}
    \vspace{-5.0pt}
    \resumeSubHeadingListStart
            \resumeProjectHeading
                {\textbf{Network Support For Scalable And High Performance Cloud Exchanges} $|$
                 \emph{\href{https://dl.acm.org/doi/abs/10.1145/3718958.3750530}{\underline{ACM SIGCOMM '25}}}}
                {Sept 2025}
            \resumeItemListStart
                % \resumeItem{\emph{M. Haseeb, J. Geng, \textbf{D. Duclos-Cavalcanti}, X. Hao, U. Butler, R. Mittal, S. Narayana, A. Sivaraman}}
                % \resumeItem{\textbf{Onyx}: Novel Cloud financial exchange achieving low latency of $<=$ 250 µs, with a difference $<$ 1 µs across 1K receivers.}
                \resumeItem{\textbf{Onyx}: cloud financial exchange delivering $\leq$\textbf{250 µs} multicast latency, within $<$ 1 µs across 1K receivers.}
                \resumeItem{\textbf{Kernel-bypass} datapath (\textbf{DPDK}) sustaining \textbf{35K packets/sec} at $\sim$\textbf{50\% lower latency} than AWS multicast.}
                % \resumeItem{Latency-minimizing \textbf{overlay multicast tree} fanning out market data, reused bidirectionally for the inbound order path.}
                % \resumeItem{Achieves better scalability and around 50\% lower latency than the multicast service provided by AWS. }
                %\resumeItem{Implemented kernel-bypass techniques (DPDK) to scale performance up to a 35K multicast packet rate.}
                %\resumeItem{Used kernel-bypass techniques (DPDK) to scale and achieve $\sim$50\% lower latency than AWS's multicast service.}
                % \resumeItem{High-performance using kernel bypassing (DPDK) to reach a throughput of 35K multicast packet rate per second.}
          \resumeItemListEnd
    \resumeSubHeadingListEnd

\vspace{-13.0pt}
\section{Projects}
    \vspace{-5.0pt}
    \resumeSubHeadingListStart
        \resumeProjectHeading
            {\href{https://github.com/duclos-cavalcanti/TreeBuilder}{\textbf{Cloud-TreeBuilder} (\emph{M.Sc. Thesis})} $|$ \emph{GCP, ZMQ, Terraform, Python, C++, Distributed Systems}}{\small{Mar 2024 -- Oct 2024}}
            \resumeItemListStart
              %\resumeItem{Launches and selects K out of N VMs in a cluster to create an optimal multicast tree of depth D and fan-out F.}
              %\resumeItem{Deploys UDP based probe jobs on VM subsets, collecting data regarding individual network performance (JSON).}
              %\resumeItem{Applies a developed heuristic on collected data to select VMs for a tree layer by layer. }
              %\resumeItem{Uses UDP probe jobs to examine VMs network performance, selecting instances for a tree layer by layer. }
              %\resumeItem{Uses terraform to manage cloud state, ZMQ for node communication and Protobufs for data serialization.}
              %\resumeItem{Leverages \href{https://www.clockwork.io/clock-sync/}{clockwork's high-precision software clock synchronization daemon} to provide a global common clock.}

            % \resumeItem{Optimally selects VMs in a cluster to form a multicast tree of depth D and fan-out F, to \textbf{minimize network latency}.}
            % \resumeItem{Deployed \textbf{UDP} probe jobs on VMs, gathering network performance data for informed heuristic selection (JSON).}
            \resumeItem{Mitigates \textbf{straggler VMs} by probing cluster network performance to build a multicast tree of depth D, fan-out F.}
            \resumeItem{Ranks candidate VMs by weighted \textbf{p90 one-way delay} and standard deviation from UDP probes.}
            \resumeItem{Developed a heuristic matching \textbf{LemonDrop}'s \textbf{NP-hard QAP} optimizer on a 25-VM cluster: \textbf{7.7\% lower p90} and \textbf{10.6\% lower median} worst-receiver latency at a fraction of the compute.}
            % \resumeItem{Integrated Terraform for cloud state management, ZMQ for node communication, and Scipy/Numpy for data analysis.}
            \resumeItemListEnd

        % \vspace{-16pt}
        % \resumeProjectHeading
        %     {\href{https://github.com/duclos-cavalcanti/rust-pkcs11}{\textbf{Rust PKCS11 Client/Server}} $|$ \emph{Rust, Multithreading, HSM, Cryptography, Protobufs}}{\small{Dec 2024 -- Present}}
        %     \resumeItemListStart
        %       \resumeItem{Implemented a PKCS11 \textbf{server} that accepts concurrent requests (\textbf{clients}) for typical cryptographic operations.}
        %       \resumeItem{Designed performant server backend with thread-safe access to SoftHSM for secure encryption, decryption and signing.}
        %     \resumeItemListEnd

        % \vspace{-16pt}
        % \resumeProjectHeading
        %     {\href{https://github.com/duclos-cavalcanti/Open-MPI-ValueIteration}{\textbf{Open-MPI Value Iteration}} $|$ \emph{C++, Parallel-Computing, MPI, HPC}}{\small{Mar 2022}}
        %     \resumeItemListStart
        %       %\resumeItem{Uses MPI techniques to iteratively distribute workload across an HPC cluster and gather results.}
        %       %\resumeItem{Uses MPI techniques to distribute workload across an HPC cluster to solve a stochastic navigation problem.}
        %       \resumeItem{An HPC prototype that solves a stochastic navigation problem by distributing workload across an HPC cluster.}
        %       \resumeItem{Leveraged MPI to scale data distribution, achieving a 52\% performance improvement over single-threaded execution.}
        %
        %       %\resumeItem{An HPC prototype that solves a stochastic navigation problem through Asynchronous Value Iteration (AVI).}
        %       %\resumeItem{Leveraged MPI to iteratively distribute workload across an HPC cluster, executing 52\% faster than in single-threaded.}
        %     \resumeItemListEnd

        % \resumeProjectHeading
        %     {\href{https://github.com/duclos-cavalcanti/microsemi-error-detection}{\textbf{Hamming Code Error Detection (16,11)}} $|$ \emph{C, VHDL, FPGA, SoC, UART}}{\small{Feb 2021}}
        %     \resumeItemListStart
        %       \resumeItem{Implemented an error detection/correction algorithm for packet transmission on Microsemi's SF2 FPGA/SoC.}
        %       \resumeItem{Error-injected packets sent between host and SoC via UART and offloaded to the FPGA for detection/correction.}
        %     \resumeItemListEnd
        %\vspace{-16pt}
        %\resumeProjectHeading
        %{\href{https://github.com/duclos-cavalcanti/FreeRTOS-SpaceInvaders}{\textbf{FreeRTOS-SpaceInvaders}} $|$ \emph{C, RTOS, Multi-Threaded}}{\small{Sept 2020}}
        %\resumeItemListStart
        %  \resumeItem{Implemented the famous arcade game as a multi-threaded FreeRTOS application in C.}
        %\resumeItemListEnd

    \resumeSubHeadingListEnd


% -------------------- SKILLS --------------------
\vspace{-10.0pt}
\section{Technical Skills}
 \begin{itemize}[leftmargin=0.15in, label={}]
    \small{\item{
    \textbf{Languages}{: Rust, C++, C, Python, Golang, Bash, VHDL} \\
    \textbf{Cloud}{: AWS (CDK, Lambda, Step Functions, DynamoDB, KMS, EventBridge), GCP, Terraform, Packer, Docker} \\
    \textbf{Tools}{: Unix Shell, Git, GitLab CI, CMake, GNU Make, Vim, Tmux} \\
    \textbf{Technologies}{: Linux, TCP/UDP, HSM, gRPC, Protobufs, ZeroMQ, Tokio, MPI, OpenMP, DPDK, FreeRTOS, FPGA} \\
    % \textbf{Languages}{: Rust, C++, C, Python, Golang, Bash, VHDL} \\
    % \textbf{Cloud Services}{: Amazon EC2 (AWS), Google Cloud Platform (GCP), Terraform, Docker, Packer, Vagrant} \\
    % \textbf{Tools}{: Linux, Unix Shell, Git, Github CI/CD, CMake, GNU Make, Bazel, Vim, GDB} \\
    % \textbf{Technologies}{: Cloud Computing, Computer Networking, Operating Systems, Machine Learning, HPC, FPGAs} \\
    % \textbf{Protocols}{: TCP, UDP, IP, ETHERNET, USB, UART} \\
    % \textbf{Hardware}{: Raspberry PIs, Embedded Linux, ARM Cortex MCUs, USB, TCP, UDP, IP, UART, GPIO} \\
    % \textbf{Verbal/Written}{: German -- Fluent, Portuguese -- Fluent}
    }}
 \end{itemize}

\vspace{-8.0pt}
%-----------EDUCATION-----------
\section{Education}
  \resumeSubHeadingListStart
    %\resumeSubheading
    \resumeSubheading
        {Technical University of Munich}{Oct 2020 -- Oct 2024}
        {M.Sc. Electrical and Computer Engineering}{Munich, Germany}
        % \resumeItemListStart
        %     % \resumeItem{Visiting Non-Degree Graduate Student: \textbf{New York University} (2023 -- 2024)}
        %     \resumeItem{Visiting Non-Degree Graduate Student: \textbf{New York University} (2023 -- 2024) -- \textit{see publications}}
        %     % \resumeItem{M.Sc. Thesis: \textbf{Virtual Machine Selection Heuristic for Multicast Overlay Trees in the Cloud}}
        %     % %\resumeItem{Co-Authored Publication: \href{https://arxiv.org/abs/2402.09527}{Design and Implementation of A Scalable Financial Exchange in the Cloud}}
        %     % \resumeItem{M.Sc. Thesis: \textbf{"Virtual Machine Selection Heuristic for Multicast Overlay Trees in the Cloud"} -- (\href{https://github.com/duclos-cavalcanti/TreeBuilder}{\textit{Link}})}
        %     %\resumeItem{M.Sc. Thesis: \textbf{"Virtual Machine Selection Heuristic for Financial Exchanges in the Cloud"} -- (\href{https://github.com/duclos-cavalcanti/TreeBuilder}{\textit{Cloud TreeBuilder}})}
        %     % \resumeItem{\textbf{Related Coursework}: Operating Systems, Machine Learning Methods, High-Performance Computing Lab}
        %     %\resumeItem{\textbf{Related Coursework}: Machine Learning Methods, HW-SW Co-Design, High Performance Computing Lab}
        % \resumeItemListEnd
    \resumeSubheading
        {Technical University of Munich}{Oct 2016 -- Sept 2020}
        {B.Sc. Electrical and Computer Engineering}{Munich, Germany}
  \resumeSubHeadingListEnd

%
\end{document}
